\subsection{Segmentering}

Segmenteringen tar utgangspunkt i tre variabler som til sammen avgjør hvilke
kundegrupper som er både teknisk tilgjengelige, kommersielt attraktive og
realistiske å konvertere.

\subsubsection{V1 -- Geografisk-teknisk tilgjengelighet}

Formålet er å identifisere hvilke kundesegmenter Ziplines plattform kan
betjene, og dermed avgrense det reelle markedet fra det teoretiske. For
dronelevering er dette en forutsetningsvariabel fremfor en
prioriteringsvariabel. Et segment med høy kjøpekraft og sterk motivasjon er
irrelevant dersom de tekniske leveringsforholdene ikke er til stede.
Variabelen operasjonaliseres gjennom boligtype, tomtestørrelse og avstand til
eksisterende matleveringsdekning. Eneboliger og småhus med uteareal på minimum
én meter i diameter er det tekniske minstekravet for Zipline, mens avstand til
nærmeste Wolt- eller Foodora-sone brukes som indikator på udekket
etterspørsel. Geografi fungerer i dette tilfellet som en indikator for flere
relevante forhold samtidig. Dersom man bor i et villastrøk i Bærum, impliserer
det med høy sannsynlighet egnet boligtype og topografiske forhold som
forsterker dronens hastighetsfordel over tradisjonell billevering.

\subsubsection{V2 -- Verdipotensial}

Formålet er å identifisere og prioritere de kundesegmentene som forventes å
generere høyest lønnsomhet over tid, basert på estimert livstidsverdi fremfor
enkeltstående bestillinger \parencite[11]{pfeifer2005}. For Zipline måles
dette ideelt sett gjennom ordrefrekvens, ordreverdi og betalingsvilje. Dersom
fullstendig data ikke er tilgjengelig, benyttes inntektsnivå og eksisterende
forbruksmønster på matleveringstjenester som alternative indikatorer. Det
legges til grunn en hypotese om at husholdninger med høy medianinntekt,
kombinert med tidspresset livssituasjon, vil ha høyere ordrefrekvens og lavere
prissensitivitet enn gjennomsnittsforbrukeren. I segmenter der Zipline er det
eneste reelle leveringsalternativet, vil betalingsviljen sannsynligvis
overstige den som observeres i markeder med konkurrerende tjenester.

\subsubsection{V3 -- Adopsjonsbarriere overfor ny leveringsteknologi}

Formålet er å identifisere hvilke kundesegmenter det er realistisk å
konvertere til Ziplines app-baserte direktekanal, gitt at dronelevering
fortsatt er et nytt konsept for norske forbrukere \parencite{rogers2003}.
Variabelen skiller mellom tre grupper. Den første består av brukere med lav
adopsjonsbarriere, som allerede er aktive på Wolt eller Foodora og åpne for ny
leveringsteknologi. Den andre består av brukere med høy barriere, som er lite
digitalt aktive og skeptiske til ubemannede systemer. Den tredje er et
påvirkelig mellomsjikt. Denne identifikasjonen er kritisk ettersom målet ikke
er å nå alle teknisk tilgjengelige husholdninger, men å prioritere de
segmentene der kommunikasjons- og lanseringstiltak faktisk kan endre atferd.
Variabelen operasjonaliseres gjennom alder, eksisterende digital aktivitet og
kjennskap til matleveringstjenester.

\subsubsection{Segment 1: Velstående villaforsteder}

Dette segmentet scorer høyt på alle tre segmenteringsvariabler og er derfor
det mest attraktive primærmarkedet for Zipline i Norge. Den geografisk-%
tekniske tilgjengeligheten er svært god. Store tomter på 500 til 2000
kvadratmeter i områder som Lommedalen, Hosle, Jar og Nesøya gir gode fysiske
forutsetninger for dronelevering. Verdipotensialet er tilsvarende høyt,
ettersom husholdningene befinner seg blant Norges rikeste delområder målt ved
medianinntekt. Adopsjonsbarrieren vurderes som lav, ettersom målgruppen består
av tidspressede toinntektsfamilier i alderen 30--50 år med urbane matvaner og
høy digital aktivitet. De er allerede vant til appbasert matlevering via Wolt
og Foodora, men har begrenset faktisk tilbud i sitt nærområde.

Verdiforslaget er særlig sterkt fordi luftlinjeavstanden til
restaurantklyngene i Sandvika og Lysaker er 5 til 10 kilometer, mens bilveien
gjennom E18 og E16 i rushtiden gjerne tar 30 til 45 minutter. Den
regulatoriske situasjonen forsterker argumentet ytterligere, ettersom U-space
sandboxen i Bærum allerede er operativ og gir Zipline et forsprang som ingen
annen norsk lokasjon kan tilby på tilsvarende tidspunkt. Segmentet vurderes
som høyeste prioritet for B2C-lansering.

\subsubsection{Segment 2: Geografisk isolerte}

Nesodden skiller seg ut som et område med høy teknisk og geografisk egnethet
for dronelevering, samtidig som dagens matleveringstilbud er fraværende. Dette
segmentet er geografisk-teknisk velegnet, med høy andel eneboliger og romslige
utearealer som legger til rette for dronelevering. Det som skiller Nesodden
fra de øvrige segmentene er imidlertid ikke graden av egnethet, men fraværet
av andre alternativer. Kommunens 21\,000 innbyggere lever med urbane
forbruksvaner og inntekter over landsgjennomsnittet, men er fysisk adskilt fra
Oslo. Veiforbindelsen er 45 kilometer lang gjennom Drøbak, mens
luftlinjeavstanden til Lysaker er knappe 5 til 8 kilometer. Verken Wolt eller
Foodora dekker kommunen i dag, og tradisjonell billevering er kommersielt lite
gunstig.

Verdipotensialet er høyt gitt inntektsnivået og den udekte etterspørselen, og
den tilnærmet monopollignende posisjonen gir grunnlag for premiumprising.
Adopsjonsbarrieren vurderes som moderat lav. Befolkningen har en urban og
pendlerpreget profil med høy digital aktivitet og god kjennskap til
matleveringstjenester, til tross for manglende tilgang i dag. Segmentet egner
seg godt for en tidlig lansering med høy kommunikasjonsverdi, og er strategisk
viktig for å demonstrere droneleveringens unike egenskaper overfor både
forbrukere og regulatorer.

\subsubsection{Segment 3: Eldre velstående}

Den geografisk-tekniske tilgjengeligheten er god i dette segmentet.
Eneboligandelen er på opptil 27 prosent, og terrenget forsterker dronens
relative fordel fremfor å hemme den. Holmenkollen og Slemdal ligger 300 til
500 meter over havet med bratte veier som gjør billevering fra
lavlandsrestauranter tidkrevende. Dronen ignorerer topografi fullstendig og
leverer på 5 til 8 minutter der biler bruker 20 til 30 minutter.
Verdipotensialet er tilsvarende høyt. Vestre Aker og Ullern har Oslos høyeste
gjennomsnittsinntekt, og Ziplines erfaringer fra Dallas--Fort Worth bekrefter
at seniorer og personer med redusert mobilitet utgjør en kjernegruppe fordi
dronelevering eliminerer behovet for bil eller fysisk mobilitet
\parencite{dronexl2026}.

Adopsjonsbarrieren er den viktigste differensiatoren for dette segmentet.
Aldersgruppen over 67 år er historisk sett mer forsiktig med å ta i bruk nye
digitale tjenester \parencite{ssb2023}, noe som innebærer at konvertering
krever en mer tilpasset innføringsstrategi enn for de øvrige segmentene.
Segmentet vurderes likevel som verdifullt og bør prioriteres som
sekundærmarked, grunnet høy betalingsvilje og stort behov.
